\section{Kết luận}

Việc cài đặt độc lập làm rõ ba cơ chế dừng. LPA yêu cầu toàn bộ nhãn cực đại sau một lượt quét. Retention-LPA khai thác một lượt không đổi. FLPA chỉ đánh giá lại các đỉnh mà lân cận có thể làm giảm ưu thế của nhãn đang giữ. Kiểm thử trên graph điển hình và thực nghiệm kiểm được maximality của những lượt hoàn thành. Giao thức chung cho phép đối chiếu thời gian và lượt xử lý thay vì suy đoán lợi ích chỉ từ tên thuật toán.

E1 ghi nhận lợi ích của lịch FIFO so với quét toàn bộ LPA và đồng thời cho thấy Retention-LPA có thể dừng nhanh khi tạo phân hoạch mịn. E3 cho thấy đổi trật tự xử lý làm thay đổi cả công việc lẫn số cụm, đặc biệt khi SBM khó phục hồi. Chúng tôi công bố seed, trạng thái dừng và dữ liệu thô để có thể kiểm những kết luận này. Tích hợp E2 và E4 sẽ bổ sung kiểm chứng chất lượng tại quy mô Figure~4 và khả năng khái quát sang liên lạc thực tế.
